
We investigated which data curation techniques transfer robustly across foundation model training stages versus which require stage-specific tuning. Through controlled experiments across C4 pre-training and Dolly/Alpaca fine-tuning using Llama-2-7B, we validated a transfer stability taxonomy categorizing techniques by objective-dependence. Low-level quality filters (deduplication, perplexity) transfer with 0.38\% average delta (95\% CI [0.30\%, 0.46\%]) and identical optimal thresholds (dedup=0.7, perplexity=500) across stages. High-level strategies (instruction quality filters) show 5.04\% degradation (95\% CI [4.44\%, 5.65\%]) when transferred, with large effect size (Cohen's d=10.76) confirming categorical separation. We quantified quality-speed trade-offs for embedding-based subset selection: early-stage embeddings achieve 98\% of late-stage quality at 15\% of compute cost (6.$9\times$ speedup, 2.4\% quality penalty).

These findings enable practitioners to reuse C4 pre-training thresholds (dedup 0.7-0.8, perplexity 500-1000) for instruction fine-tuning without re-optimization, focusing tuning resources on domain mixing and task-specific filters. This shifts the paradigm from per-stage curation optimization to transfer-aware design.

\textbf{Methodological constraint:} All experiments conducted at proof-of-concept level. Curation pipelines executed on full datasets, but model training simulated using predicted scores from published benchmarks (Llama-2 technical report, DataComp). This validates mechanism direction---categorical separation, threshold invariance---but defers absolute precision to production execution. Production validation with full training and lm-evaluation-harness needed to confirm absolute performance values.

\subsubsection{Future Directions}

\textbf{Untested scope extensions:} Our validation tested text-only models across pre-training $\rightarrow$ fine-tuning. The taxonomy predicts image deduplication (pHash) should transfer robustly across vision-language pre-training and fine-tuning, while aesthetic scoring (objective-dependent) requires stage-specific tuning. For RLHF, safety filters (universal hygiene) should transfer from pre-training while preference alignment criteria (objective-dependent) require RLHF-specific optimization. These predictions await empirical validation.

\textbf{Conservative experimental settings:} Our test datasets exhibited low duplicate burden (0.03-0.10\%), yielding conservative curation effects. High-noise production datasets (web-scraped fine-tuning corpora) would show larger absolute impact while preserving transfer patterns. Testing on such data would validate categorical separation under higher curation pressure.

\textbf{Model scale assumptions:} We validated at 7B parameter scale (Llama-2-7B). The universal hygiene hypothesis predicts transfer robustness should hold across model scales (13B, 70B, 175B)---optimal deduplication thresholds encode data quality independent of model capacity. Replicating h-m1 threshold transfer at multiple scales would confirm scale-invariance.

\textbf{Boundary condition testing:} We tested moderate distribution shift (C4 web text $\rightarrow$ Dolly instructions). The taxonomy predicts transfer delta increases with domain distance. High-shift scenarios (biomedical literature fine-tuning, legal document pre-training) would define breakdown thresholds where even low-level filters require re-tuning. Characterizing the transfer delta = f(domain\_distance) curve would refine applicability boundaries.

As foundation models continue to scale across training stages---from trillion-token pre-training to specialized fine-tuning to RLHF---understanding which data curation decisions transfer will become increasingly critical for efficient development. Our taxonomy provides a first step toward transfer-aware curation design, with production validation needed for trillion-token scale and RLHF stages.
